\documentclass[11pt,letterpaper]{article}

% ==============================================================================
%% Encoding & idioma
% ==============================================================================
\usepackage[utf8]{inputenc}
\usepackage[T1]{fontenc}
\usepackage[spanish]{babel}

% ==============================================================================
%% Tipografía y diseño
% ==============================================================================
\usepackage{lmodern}
\usepackage[margin=2.5cm]{geometry}
\usepackage{xcolor}
\usepackage{enumitem}
\usepackage{titlesec}
\usepackage{parskip}
\usepackage{array}
\usepackage{booktabs}
\usepackage{hyperref}

% ==============================================================================
%% Paleta de colores (consistente con plantilla Beamer Colmex)
% ==============================================================================
\definecolor{main}{HTML}{5E002B}
\definecolor{subtitle}{RGB}{75,75,75}
\definecolor{cerulean}{RGB}{0,123,167}
\definecolor{redorange}{RGB}{210,77,49}
\definecolor{forestgreen}{RGB}{36,128,103}

\hypersetup{
    colorlinks=true,
    linkcolor=main,
    citecolor=main,
    urlcolor=cerulean,
    pdftitle={Programación para Proyectos de Datos -- Syllabus},
    pdfauthor={Daniel Kelly}
}

% ==============================================================================
%% Estilo de secciones
% ==============================================================================
\titleformat{\section}
  {\Large\bfseries\color{main}}{\thesection}{1em}{}
\titleformat{\subsection}
  {\large\bfseries\color{main}}{\thesubsection}{1em}{}
\titleformat{\subsubsection}
  {\normalsize\bfseries\color{subtitle}}{}{0em}{}

% ==============================================================================
%% Comandos auxiliares (consistentes con plantilla Beamer)
% ==============================================================================
\newcommand{\blue}[1]{{\color{cerulean}#1}}
\newcommand{\orange}[1]{{\color{redorange}#1}}
\newcommand{\green}[1]{{\color{forestgreen}#1}}

% Marcador de contenido pendiente (para iterar sesión por sesión)
\newcommand{\porDesarrollar}{\textit{\color{redorange}[Por desarrollar]}}

% Divisor de parte del curso: nueva página, título grande, espacio respiratorio.
\newcommand{\parteCurso}[1]{%
  \clearpage%
  \vspace*{0.5cm}%
  {\centering\color{main}\fontsize{26pt}{32pt}\selectfont\bfseries #1\par}%
  \vspace{0.8cm}%
}

% Bloque estandarizado para cada semana. No fuerza \clearpage: las semanas
% fluyen continuas dentro de cada parte. El título se renderiza prominente
% (\Large bold, color vino) para mantener la jerarquía visual.
\newcommand{\semana}[2]{%
  \bigskip%
  {\color{main}\Large\bfseries Semana #1: #2\par}%
  \smallskip%
}

% Sub-bloque para cada sesión dentro de una semana
\newcommand{\subsesion}[2]{%
  \medskip\noindent\textbf{\color{main}Sesión #1 --- #2}\par\smallskip%
}

% Referencia a cheatsheets oficiales de Posit (rstudio.github.io/cheatsheets)
% Uso: \cheatsheet{slug-del-archivo}{Título descriptivo}
\newcommand{\cheatsheet}[2]{%
  \href{https://rstudio.github.io/cheatsheets/#1.pdf}{\textit{Cheatsheet: #2}}%
}

% ==============================================================================
%% Bibliografía
% ==============================================================================
\usepackage[style=apa, backend=biber]{biblatex}
\addbibresource{../../bibliography.bib}

% ==============================================================================
%% Encabezado
% ==============================================================================
\title{\color{main}\textbf{Programación para Proyectos de Datos}\\[0.3em]
       \large\color{subtitle}Temario del Curso}
\author{Daniel Kelly\\
        \small El Colegio de México --- Licenciatura en Economía}
\date{Otoño 2026}

\begin{document}
\maketitle

% ==============================================================================
\section{Presentación}
% ==============================================================================

El curso aborda la construcción de proyectos de datos como flujo integrado, no como yuxtaposición de habilidades. Los cursos previos de la Licenciatura en Economía de El Colegio de México proveen al estudiante de competencias específicas en programación y análisis estadístico; el énfasis aquí está en su articulación dentro de un mismo repositorio: extracción de datos de fuentes heterogéneas, limpieza sistemática, modelado y producción de un entregable comunicable, bajo prácticas reproducibles.

El curso se imparte bajo la filosofía \textit{zero-to-hero}: se asume que el estudiante puede llegar sin experiencia previa de programación y se espera que cierre el semestre capaz de levantar un proyecto de datos completo por su cuenta.

El lenguaje del curso es R, pensado como una habilidad extremadamente relevante en el campo de las finanzas, la economía y el análisis de políticas públicas. Como IDE (\textit{Integrated Development Environment}) se recomienda usar Positron, el hermano más moderno de R-Studio. La elección de R sobre alternativas (Python, Stata) se discute formalmente en la primera sesión.

% ==============================================================================
\section{Objetivos}
% ==============================================================================

Al finalizar el curso, el estudiante será capaz de:

\begin{enumerate}[leftmargin=*]
    \item Estructurar un proyecto de datos completo bajo convenciones reproducibles: control de versiones, organización de directorios y gestión de dependencias.
    \item Importar, limpiar y transformar datos de fuentes heterogéneas: archivos planos, hojas de cálculo, APIs HTTP y bases de datos relacionales.
    \item Programar funciones reutilizables y aplicar paradigmas de programación funcional para automatizar tareas repetitivas.
    \item Estimar y reportar modelos estadísticos comunes en investigación económica desde una perspectiva programática.
    \item Construir entregables comunicables: visualizaciones, documentos automatizados, aplicaciones interactivas e integraciones con servicios externos.
\end{enumerate}

% ==============================================================================
\section{Prerrequisitos}
% ==============================================================================

\begin{itemize}[leftmargin=*]
    \item No se requiere experiencia previa en programación.
    \item Se asume manejo básico de estadística descriptiva y nociones iniciales de econometría, incluyendo regresión lineal univariable.
    \item El curso no presenta la teoría estadística subyacente a los modelos; solo su implementación programática.
\end{itemize}

% ==============================================================================
\section{Metodología}
% ==============================================================================

\textbf{Duración.} 16 semanas, con 2 sesiones de 1:30 hr por semana (32 sesiones totales).

\textbf{Estructura semanal.} Cada semana tiene dos sesiones con funciones diferenciadas:

\begin{itemize}[leftmargin=*]
    \item \textbf{\blue{Sesión 1 --- Teoría.}} Exposición conceptual del tema, con código en vivo como ilustración. Proporción típica: 60\% teoría, 40\% código demostrativo.
    \item \textbf{\orange{Sesión 2 --- Práctica.}} Resumen breve de la sesión 1 y ejercicios resueltos en clase.
\end{itemize}

\textbf{Tareas.} El curso no contempla tareas semanales fuera de clase. El trabajo evaluable se concentra en los ejercicios de la segunda sesión y en el proyecto integrador.

% ==============================================================================
\section{Evaluación}
% ==============================================================================

\begin{center}
\begin{tabular}{lr}
\toprule
\textbf{Componente} & \textbf{Peso} \\
\midrule
Checkpoints del proyecto integrador           & 50\% \\
Proyecto final                                & 20\% \\
Ejercicios en clase (segunda sesión semanal)  & 20\% \\
Asistencia                                    & 10\%  \\
\bottomrule
\end{tabular}
\end{center}

% ==============================================================================
\section{Estructura del curso}
% ==============================================================================

El curso se organiza en cinco partes a lo largo de 16 semanas:

\smallskip
\noindent\textbf{\textcolor{main}{Parte I --- Fundamentos}}
\begin{itemize}[leftmargin=*, itemsep=1pt]
    \item \textbf{Semana 1.} ¿Por qué R? Setup y primeros pasos.
    \item \textbf{Semana 2.} Estructuras de datos e indexación.
    \item \textbf{Semana 3.} Importación de datos y EDA con base R.
    \item \textbf{Semana 4.} Estructura de proyectos, Git y reproducibilidad.
\end{itemize}

\smallskip
\noindent\textbf{\textcolor{main}{Parte II --- Manipulación}}
\begin{itemize}[leftmargin=*, itemsep=1pt]
    \item \textbf{Semana 5.} \texttt{dplyr} I: \textit{pipe}, verbos básicos y sistema tidyverse.
    \item \textbf{Semana 6.} \texttt{dplyr} II: \textit{joins}, \textit{pivots} y manejo de \texttt{NA}.
    \item \textbf{Semana 7.} Strings, expresiones regulares y \texttt{stringr}.
\end{itemize}

\smallskip
\noindent\textbf{\textcolor{main}{Parte III --- Programación}}
\begin{itemize}[leftmargin=*, itemsep=1pt]
    \item \textbf{Semana 8.} Control de flujo y vectorización.
    \item \textbf{Semana 9.} Funciones: diseño, \textit{scoping} y \textit{debugging}.
    \item \textbf{Semana 10.} Programación funcional: \texttt{purrr} y la familia \texttt{apply}.
\end{itemize}

\smallskip
\noindent\textbf{\textcolor{main}{Parte IV --- Análisis}}
\begin{itemize}[leftmargin=*, itemsep=1pt]
    \item \textbf{Semana 11.} Modelado estadístico programático.
    \item \textbf{Semana 12.} Datos a escala: SQL, \texttt{DBI} y \texttt{dbplyr}.
\end{itemize}

\smallskip
\noindent\textbf{\textcolor{main}{Parte V --- Productos Finales}}
\begin{itemize}[leftmargin=*, itemsep=1pt]
    \item \textbf{Semana 13.} Visualización con \texttt{ggplot2}.
    \item \textbf{Semana 14.} Automatización con \texttt{officer}, \texttt{googledrive} y \texttt{gmailr}.
    \item \textbf{Semana 15.} Aplicaciones interactivas con \texttt{Shiny}.
    \item \textbf{Semana 16.} Interfaces de modelos de lenguaje.
\end{itemize}

% ------------------------------------------------------------------------------
\parteCurso{Parte I --- Fundamentos}
% ------------------------------------------------------------------------------

\semana{1}{¿Por qué R? Setup y primeros pasos}

Sesión de apertura. Cubre la justificación del lenguaje, la instalación y configuración del entorno de trabajo, los conceptos iniciales del lenguaje y la política del curso sobre el uso de modelos de lenguaje.

\textbf{Subtemas:}

\begin{itemize}[leftmargin=*, itemsep=2pt]
    \item \textbf{Bienvenida y motivación del curso.} Estructura del curso y esquema de evaluación.
    \item \textbf{¿Por qué R?} R en el contexto de la investigación económica académica; contraste con Python; razones para descartar Stata. Ecosistema CRAN, tidyverse y Posit.
    \item \textbf{Setup: Positron + R.} Instalación; layout del IDE (consola, source, environment, plots, packages); configuración mínima.
    \item \textbf{Uso responsable de modelos de lenguaje.} Distinción operativa entre el uso de un LLM como \blue{chatbot} (pregunta y respuesta aisladas) y como \orange{agente integrado al proyecto} (lectura de archivos, propuesta de cambios y ejecución de comandos con supervisión humana en el ciclo de revisión), modalidad propia de herramientas como Claude Code de Anthropic o Codex de OpenAI. \textbf{Política del curso:} los agentes están permitidos en el proyecto integrador y no permitidos en los ejercicios de clase. El estudiante debe ser capaz de leer y modificar cualquier línea que entregue como propia.
    \item \textbf{Características del lenguaje.} R como lenguaje interpretado, vectorizado y multiparadigma. Tres modos de ejecución: REPL, \textit{script} con \texttt{source()}, ejecutable con \texttt{Rscript}.
    \item \textbf{Primeros objetos en R.} Asignación (\texttt{<-} frente a \texttt{=}); tipos atómicos (\texttt{numeric}, \texttt{integer}, \texttt{logical}, \texttt{character}); construcción de vectores con \texttt{c()}, vectorización aritmética; funciones de inspección (\texttt{class}, \texttt{typeof}, \texttt{length}, \texttt{str}).
    \item \textbf{Evaluación de condiciones lógicas.} Operadores de comparación (\texttt{==}, \texttt{!=}, \texttt{<}, \texttt{>}, \texttt{<=}, \texttt{>=}) y operadores lógicos vectorizados (\texttt{\&}, \texttt{|}, \texttt{!}) frente a sus contrapartes escalares (\texttt{\&\&}, \texttt{||}). Coerción de \texttt{logical} a numérico: \texttt{sum()} sobre un vector lógico cuenta los \texttt{TRUE}; \texttt{mean()} calcula la proporción. Funciones auxiliares (\texttt{any}, \texttt{all}, \texttt{xor}, \texttt{isTRUE}, \texttt{isFALSE}) y propagación de \texttt{NA} en condiciones.
\end{itemize}

\medskip
\noindent\textbf{Objetivos de aprendizaje}
\begin{enumerate}[leftmargin=*, itemsep=1pt]
    \item Justificar la elección de R sobre alternativas programáticas en el contexto de un proyecto académico de investigación económica.
    \item Configurar Positron + R y ejecutar un \textit{script} básico tanto desde el IDE como desde la terminal.
    \item Crear, asignar y operar con vectores atómicos, identificando su tipo y tamaño.
    \item Distinguir el uso de un LLM como chatbot frente a su uso como agente integrado al proyecto, y la política del curso al respecto.
\end{enumerate}

\noindent\textbf{Lecturas}
\begin{itemize}[leftmargin=*, itemsep=1pt]
    \item \textit{R for Data Science} (2.\textsuperscript{a} ed.), \textit{Introduction} y Cap.\ 2 \textit{Workflow: basics} \cite{wickham2023r4ds}.
    \item \textit{Advanced R} (2.\textsuperscript{a} ed.), Cap.\ 2 §2.1--2.2 y Cap.\ 3 §3.1--3.2 \cite{wickham2019advr}.
\end{itemize}

\semana{2}{Estructuras de datos e indexación}

Extensión del vocabulario del lenguaje más allá de los vectores atómicos. La semana introduce listas, \texttt{data.frame} y \texttt{tibble} como las estructuras de trabajo del resto del semestre, junto con dos mecanismos transversales: la indexación de elementos y la coerción implícita de tipos. Cierra con el tratamiento de valores faltantes.

\textbf{Subtemas:}

\begin{itemize}[leftmargin=*, itemsep=2pt]
    \item \textbf{Listas.} Contenedores heterogéneos; nombres; acceso con \texttt{\$} y con \texttt{[[]]}; diferencia operativa con vectores atómicos. Uso típico: agrupar resultados de tipos distintos provenientes de un mismo análisis.
    \item \textbf{DataFrames y tibbles.} Estructura interna (lista de vectores de igual longitud con clase \texttt{data.frame}); constructores \texttt{data.frame()} y \texttt{tibble()}; funciones de inspección (\texttt{names}, \texttt{nrow}, \texttt{ncol}, \texttt{dim}, \texttt{head}, \texttt{str}, \texttt{glimpse}). Diferencias prácticas entre \texttt{tibble} y \texttt{data.frame}: impresión, ausencia de \textit{partial matching}, no conversión automática de \textit{strings}, \textit{subsetting} que no simplifica. Por convención el curso utiliza \texttt{tibble} por defecto.
    \item \textbf{Indexación.} Operadores \texttt{[}, \texttt{[[} y \texttt{\$}: preservación de clase frente a extracción de elemento. Indexación posicional, por nombre, lógica y negativa. Función \texttt{which()}. Asignación por \textit{subsetting}.
    \item \textbf{Coerción de tipos.} Coerción implícita en la construcción de vectores y en operaciones aritméticas; jerarquía \texttt{logical} $\to$ \texttt{integer} $\to$ \texttt{double} $\to$ \texttt{character}. Coerción explícita con \texttt{as.numeric()}, \texttt{as.character()}, \texttt{as.logical()}. Caso recurrente en datos de encuesta: variables numéricas almacenadas como \textit{string}.
    \item \textbf{Valores faltantes.} Naturaleza del \texttt{NA}; propagación en operaciones; \texttt{is.na()}; argumento \texttt{na.rm = TRUE} en funciones agregadoras. Distinción entre \texttt{NA}, \texttt{NULL} y \texttt{NaN}, y casos en que aparece cada uno.
\end{itemize}

\medskip
\noindent\textbf{Objetivos de aprendizaje}
\begin{enumerate}[leftmargin=*, itemsep=1pt]
    \item Distinguir entre vectores atómicos, listas, DataFrames y tibbles, y elegir la estructura adecuada según los datos a representar.
    \item Aplicar los operadores \texttt{[}, \texttt{[[} y \texttt{\$} según el resultado deseado: preservación de clase frente a extracción.
    \item Anticipar las consecuencias de la coerción implícita al construir o transformar vectores con tipos mixtos.
    \item Identificar y manejar valores faltantes en operaciones, distinguiéndolos de \texttt{NULL} y \texttt{NaN}.
\end{enumerate}

\medskip
\noindent\textbf{Lecturas}
\begin{itemize}[leftmargin=*, itemsep=1pt]
    \item \textit{Advanced R} (2.\textsuperscript{a} ed.), Cap.\ 3 \textit{Vectors} (§3.2.3, §3.5, §3.6) y Cap.\ 4 \textit{Subsetting} completo \cite{wickham2019advr}.
    \item \textit{R for Data Science} (2.\textsuperscript{a} ed.), Cap.\ 18 \textit{Missing values} \cite{wickham2023r4ds}.
\end{itemize}

\medskip
\noindent\textbf{Referencia rápida:} \cheatsheet{base-r}{Base R} (Posit).

\semana{3}{Importación de datos y EDA con base R}

Primera aproximación al trabajo con datos reales. La semana cubre el flujo inicial de cualquier proyecto: importación desde archivos o APIs HTTP, inspección post-importación, estadística descriptiva y visualización exploratoria. Introduce además el lenguaje del Análisis Exploratorio de Datos (EDA) como disciplina.

El uso se restringe deliberadamente a las funciones base de R. La decisión es pedagógica: el estudiante reconoce qué ofrece el lenguaje por sí mismo antes de incorporar el tidyverse (Semanas 5--7) y \texttt{ggplot2} (Semana 13).

\textbf{Subtemas:}

\begin{itemize}[leftmargin=*, itemsep=2pt]
    \item \textbf{Importación desde archivos.} Lectura de CSV con \texttt{readr} (\texttt{read\_csv}; diferencia con \texttt{read.csv} de base; \texttt{problems()}, \texttt{col\_types}) y de Excel con \texttt{readxl} (\texttt{read\_excel}, \texttt{sheets}, \texttt{range}). Argumentos relevantes: \textit{encoding} (UTF-8 frente a Latin-1, recurrente en datos mexicanos), representaciones de \texttt{NA}, \texttt{locale}.
    \item \textbf{Consumo de APIs.} Datos como producto de servicios web. Introducción a \texttt{httr2} como cliente HTTP moderno: construcción de \textit{requests} con \texttt{request()}, ejecución con \texttt{req\_perform()}, parseo de respuestas JSON con \texttt{resp\_body\_json()}, autenticación por \textit{token} en \textit{headers}. Caso de estudio: la API SIE de Banco de México para series macroeconómicas (tipo de cambio, INPC, tasas de interés).
    \item \textbf{Inspección post-importación.} Verificación de la importación con \texttt{str}, \texttt{glimpse}, \texttt{summary}, \texttt{head}, \texttt{tail}, \texttt{dim}. Errores frecuentes: tipos mal inferidos, \texttt{NA}s no reconocidos, \textit{encoding} corrupto.
    \item \textbf{Estadística descriptiva con base R.} Univariadas: \texttt{mean}, \texttt{median}, \texttt{var}, \texttt{sd}, \texttt{quantile}, \texttt{range}, \texttt{summary}. Frecuencias: \texttt{table}, \texttt{prop.table}. Bivariadas: \texttt{cor}, \texttt{cov}, tablas cruzadas con \texttt{table()}. Manejo de \texttt{NA} en agregaciones (\texttt{na.rm = TRUE}).
    \item \textbf{Visualización exploratoria con \textit{base graphics}.} \texttt{plot}, \texttt{hist}, \texttt{boxplot}, \texttt{barplot}, \texttt{pairs}. Argumentos básicos (\texttt{main}, \texttt{xlab}, \texttt{ylab}, \texttt{col}, \texttt{pch}) y anotaciones (\texttt{abline}, \texttt{lines}, \texttt{points}, \texttt{text}). Criterios para elegir tipo de gráfica según la pregunta.
    \item \textbf{Flujo de EDA.} Iteración entre pregunta, carga, inspección, descriptivas, visualización y ajuste de pregunta. Documentación de hallazgos.
\end{itemize}

\medskip
\noindent\textbf{Objetivos de aprendizaje}
\begin{enumerate}[leftmargin=*, itemsep=1pt]
    \item Importar datos desde archivos (CSV, Excel) verificando \textit{encoding}, tipos y representación de \texttt{NA}.
    \item Consumir APIs HTTP con \texttt{httr2} y parsear respuestas JSON, usando la API SIE de Banco de México como caso de estudio.
    \item Realizar una inspección sistemática post-importación que detecte errores tempranos (tipos mal inferidos, \texttt{NA}s no reconocidos).
    \item Calcular estadísticas descriptivas univariadas y bivariadas con funciones base, manejando los \texttt{NA} adecuadamente.
    \item Producir gráficas exploratorias con \textit{base graphics}, eligiendo el tipo apropiado según la pregunta.
\end{enumerate}

\medskip
\noindent\textbf{Lecturas}
\begin{itemize}[leftmargin=*, itemsep=1pt]
    \item \textit{R for Data Science} (2.\textsuperscript{a} ed.), Cap.\ 7 \textit{Data import} y Cap.\ 20 \textit{Spreadsheets} \cite{wickham2023r4ds}.
    \item \textit{R for Data Science} (2.\textsuperscript{a} ed.), Cap.\ 10 \textit{Exploratory data analysis} (marco conceptual; los ejemplos visuales usan \texttt{ggplot2}, cubierto en la Semana 13) \cite{wickham2023r4ds}.
    \item \textit{An Introduction to R} (manual oficial), Cap.\ 12 \textit{Graphical procedures} \cite{rcoreteam-intro}.
\end{itemize}

\medskip
\noindent\textbf{Referencia rápida:} \cheatsheet{data-import}{Data import with \texttt{readr}} (Posit).

\semana{4}{Estructura de proyectos, Git y reproducibilidad}

Cierre de la Parte I. La semana presenta las herramientas que convierten un \textit{script} aislado en un proyecto reproducible: estructura de directorios, control de versiones con Git y captura de dependencias con \texttt{renv}. La estructura propuesta proviene de la convención de proyectos profesionales y se simplifica para el contexto académico.

\textbf{Subtemas:}

\begin{itemize}[leftmargin=*, itemsep=2pt]
    \item \textbf{Estructura de un proyecto de datos.} Convención del curso: \texttt{files/} (insumos crudos), \texttt{output/} (datos procesados, gráficas, reportes), \texttt{docs/} (documentación, descriptores) y \texttt{pre/} (\textit{scripts} de \textit{preprocessing} y ETL). La separación responde a criterios distintos por carpeta: versionado, tamaño y audiencia.
    \item \textbf{Proyectos en Positron (\texttt{.Rproj}).} Función del archivo de proyecto: \textit{working directory} fijado al \textit{root} y configuración aislada. Con el proyecto abierto en Positron, los \textit{paths} relativos resuelven sin herramientas adicionales y los \textit{paths} absolutos pierden sentido. Flujos autónomos (\textit{scripts} desde terminal, \textit{scheduling}) quedan fuera del alcance del curso.
    \item \textbf{Control de versiones con Git.} Comandos básicos: \texttt{init}, \texttt{add}, \texttt{commit}, \texttt{status}, \texttt{log}, \texttt{diff}, \texttt{branch}, \texttt{checkout}. Concepto de \textit{remote}; sincronización con GitHub (\texttt{push}, \texttt{pull}). Convenciones de mensajes de \textit{commit}.
    \item \textbf{\texttt{.gitignore}: archivos no versionados.} Datos crudos pesados (\texttt{files/}), credenciales (\texttt{.env}), \textit{outputs} grandes, archivos de sistema (\texttt{.DS\_Store}), historial de R (\texttt{.Rhistory}, \texttt{.RData}). Implicaciones de versionar datos crudos o credenciales.
    \item \textbf{Reproducibilidad con \texttt{renv}.} \textit{Snapshot} de las dependencias del proyecto con \texttt{renv::init()}, \texttt{renv::snapshot()}, \texttt{renv::restore()}. Criterios para incorporar \texttt{renv}: proyectos colaborativos o entregables académicos.
    \item \textbf{README y \textit{master script}.} El \texttt{README.md} como pieza central: descripción del proyecto, instrucciones de ejecución, estructura de directorios, procedimiento para regenerar los \textit{outputs}. Convención del \textit{master script} (\texttt{pre/master.R}) que orquesta el ETL completo mediante \texttt{source()}.
\end{itemize}

\medskip
\noindent\textbf{Objetivos de aprendizaje}
\begin{enumerate}[leftmargin=*, itemsep=1pt]
    \item Diseñar la estructura de directorios de un proyecto de datos según la convención del curso.
    \item Versionar un proyecto con Git, distinguiendo qué archivos se rastrean y cuáles deben excluirse mediante \texttt{.gitignore}.
    \item Utilizar \textit{paths} relativos al \textit{root} del proyecto, garantizando la ejecución correcta de los \textit{scripts} en cualquier máquina.
    \item Capturar y restaurar las dependencias del proyecto con \texttt{renv}.
\end{enumerate}

\medskip
\noindent\textbf{Lecturas}
\begin{itemize}[leftmargin=*, itemsep=1pt]
    \item \textit{R for Data Science} (2.\textsuperscript{a} ed.), Cap.\ 6 \textit{Workflow: scripts and projects} \cite{wickham2023r4ds}.
    \item \textit{Happy Git and GitHub for the useR} (Bryan): capítulos introductorios sobre \textit{setup} y operaciones básicas con Git \cite{bryan-happygit}.
    \item Documentación oficial de \texttt{renv}, secciones \textit{Introduction} y \textit{Workflow} \cite{ushey-renv}.
\end{itemize}

% ------------------------------------------------------------------------------
\parteCurso{Parte II --- Manipulación}
% ------------------------------------------------------------------------------

\semana{5}{dplyr I: pipe y verbos básicos}

Apertura de la Parte II. El contenido subsiguiente del curso ---limpieza avanzada, programación, modelado y productos--- presupone fluidez con \texttt{dplyr} y con el \textit{pipe} como mecanismo de composición. La meta de la semana es alcanzar dicha fluidez: traducir una pregunta sobre un \texttt{tibble} en una secuencia de verbos, y leer código tidyverse de terceros sin interrupciones.

\textbf{Subtemas:}

\begin{itemize}[leftmargin=*, itemsep=2pt]
    \item \textbf{El sistema tidyverse.} Paquetes que componen el ecosistema: \texttt{dplyr}, \texttt{tidyr}, \texttt{ggplot2}, \texttt{readr}, \texttt{tibble}, \texttt{purrr}, \texttt{forcats}, \texttt{stringr}, \texttt{lubridate}. Gramática y vocabulario compartidos frente a la alternativa de paquetes desacoplados. Convención del curso: carga mediante \texttt{pacman::p\_load(tidyverse)}.
    \item \textbf{Tidy data.} Tres reglas: cada variable es una columna, cada observación una fila, cada valor una celda. Diagnóstico de formatos no \textit{tidy} típicos en encuestas (variables-columna con sufijos por año o módulo). El \textit{reshape} hacia \textit{tidy} se introduce conceptualmente; el detalle operativo con \texttt{pivot\_longer()} y \texttt{pivot\_wider()} se aborda en la Semana 6.
    \item \textbf{El \textit{pipe}.} \textit{Pipe} nativo \texttt{|>} (R 4.1+) y \textit{pipe} del tidyverse \texttt{\%>\%} (\texttt{magrittr}). Semántica: el resultado del lado izquierdo se inserta como primer argumento del derecho. Convención del curso: \textit{pipe} nativo por defecto; \texttt{\%>\%} solo cuando se requieren sus capacidades específicas.
    \item \textbf{Los seis verbos de \texttt{dplyr}.} \texttt{filter()} (filas por condición), \texttt{arrange()} (ordenamiento), \texttt{select()} (columnas, con \textit{helpers} \texttt{starts\_with}, \texttt{ends\_with}, \texttt{contains}, \texttt{where}), \texttt{mutate()} (creación o modificación), \texttt{summarize()} (colapso a resumen), \texttt{group\_by()} y \texttt{ungroup()} (contexto de grupo). \textit{Helpers} frecuentes: \texttt{slice\_*()}, \texttt{count()}, \texttt{distinct()}, \texttt{if\_else()}, \texttt{between()}, \texttt{n()}, \texttt{n\_distinct()}. Los verbos como bloques componibles: toda transformación es una secuencia conectada por el \textit{pipe}.
    \item \textbf{\textit{Non-Standard Evaluation} (NSE).} Justificación de por qué \texttt{filter(df, x > 5)} funciona sin \texttt{df\$x} ni \texttt{"x"}. Introducción suficiente para anticipar las limitaciones que aparecen al escribir funciones que reciben nombres de columnas como argumentos. El tratamiento operativo (\texttt{\{\{~\}\}}, \textit{tunneling}) se difiere a la Semana 9.
\end{itemize}

\medskip
\noindent\textbf{Objetivos de aprendizaje}
\begin{enumerate}[leftmargin=*, itemsep=1pt]
    \item Cargar el ecosistema tidyverse y reconocer qué paquete provee qué funcionalidad.
    \item Identificar si un dataset está en formato \textit{tidy} y, en su caso, conceptualizar el \textit{reshape} requerido.
    \item Componer secuencias de verbos \texttt{dplyr} conectados por el \textit{pipe} para resolver preguntas de manipulación de datos.
    \item Leer código tidyverse escrito por terceros sin necesidad de ejecutarlo línea por línea.
\end{enumerate}

\medskip
\noindent\textbf{Lecturas}
\begin{itemize}[leftmargin=*, itemsep=1pt]
    \item \textit{R for Data Science} (2.\textsuperscript{a} ed.), Cap.\ 3 \textit{Data transformation} completo \cite{wickham2023r4ds}.
    \item \textit{R for Data Science} (2.\textsuperscript{a} ed.), Cap.\ 5 \textit{Data tidying} (§5.1--§5.2); §5.3--§5.4 sobre \textit{pivots} corresponden a la Semana 6 \cite{wickham2023r4ds}.
\end{itemize}

\medskip
\noindent\textbf{Referencia rápida:} \cheatsheet{data-transformation}{Data transformation with \texttt{dplyr}} (Posit).

\semana{6}{dplyr II: joins, pivots y manejo de NA}

La semana completa las herramientas necesarias para transformar datos crudos en un dataset \textit{tidy} listo para análisis. Cubre los \textit{pivots} que cierran el tema de \textit{tidy data} iniciado en la Semana 5, los \textit{joins} entre tablas, el manejo avanzado de \texttt{NA} y la construcción de variables categóricas derivadas. Al cierre, el estudiante dispone del conjunto suficiente para armar el pipeline E$\to$T del flujo planteado al inicio del curso.

\textbf{Subtemas:}

\begin{itemize}[leftmargin=*, itemsep=2pt]
    \item \textbf{Pivots.} \texttt{pivot\_longer()} y \texttt{pivot\_wider()} para \textit{reshape} entre formato ancho y largo. Argumentos: \texttt{cols}, \texttt{names\_to}, \texttt{values\_to} (\textit{longer}); \texttt{names\_from}, \texttt{values\_from} (\textit{wider}). Casos recurrentes en encuestas: variables con sufijos por año o módulo (ENSU, ENIGH).
    \item \textbf{\textit{Joins} entre tablas.} \textit{Mutating joins} (\texttt{left\_join}, \texttt{right\_join}, \texttt{inner\_join}, \texttt{full\_join}) frente a \textit{filtering joins} (\texttt{semi\_join}, \texttt{anti\_join}). Concepto de \textit{key} y sintaxis recomendada con \texttt{join\_by()}. Errores frecuentes: cardinalidad inesperada (\textit{many-to-many}), \textit{keys} duplicadas, \textit{joins} silenciosos sobre columnas mal pareadas. La semana introduce \texttt{tidylog}, que intercepta operaciones de \texttt{dplyr} y \texttt{tidyr} y emite en consola un reporte de filas afectadas, ausencias de \textit{match} y aparición de \texttt{NA} nuevos. Aplicable a todos los verbos vistos en las dos últimas semanas.
    \item \textbf{Manejo avanzado de \texttt{NA}.} \texttt{replace\_na()} para rellenar valores específicos; \texttt{coalesce()} para tomar el primer valor no \texttt{NA} entre varias columnas; \texttt{drop\_na()} para descartar filas con \texttt{NA} en columnas indicadas. La decisión de propagar o intervenir debe ser explícita y documentada.
    \item \textbf{Construcción de variables categóricas.} \texttt{if\_else()} vectorizado para condiciones binarias; \texttt{case\_when()} para reglas múltiples mutuamente excluyentes. Mención a \texttt{forcats} para reordenar y recodificar factores, sin profundizar.
\end{itemize}

\medskip
\noindent\textbf{Objetivos de aprendizaje}
\begin{enumerate}[leftmargin=*, itemsep=1pt]
    \item Realizar \textit{reshape} de un dataset entre formatos ancho y largo según la operación requerida.
    \item Combinar múltiples tablas con los \textit{joins} apropiados, identificando \textit{keys} y anticipando problemas de cardinalidad.
    \item Decidir explícitamente el manejo de \texttt{NA} durante la limpieza, distinguiendo cuándo propagarlos y cuándo intervenir.
    \item Construir variables derivadas categóricas mediante reglas explícitas con \texttt{if\_else()} y \texttt{case\_when()}.
\end{enumerate}

\medskip
\noindent\textbf{Lecturas}
\begin{itemize}[leftmargin=*, itemsep=1pt]
    \item \textit{R for Data Science} (2.\textsuperscript{a} ed.), Cap.\ 5 \textit{Data tidying}, §5.3 \textit{Lengthening data} y §5.4 \textit{Widening data} \cite{wickham2023r4ds}.
    \item \textit{R for Data Science} (2.\textsuperscript{a} ed.), Cap.\ 19 \textit{Joins} completo \cite{wickham2023r4ds}.
    \item \textit{R for Data Science} (2.\textsuperscript{a} ed.), Cap.\ 18 \textit{Missing values}: revisión operativa de \texttt{replace\_na()}, \texttt{coalesce()} y \texttt{drop\_na()} \cite{wickham2023r4ds}.
\end{itemize}

\medskip
\noindent\textbf{Referencia rápida:} \cheatsheet{tidyr}{Data tidying with \texttt{tidyr}} (Posit).

\semana{7}{Strings, expresiones regulares y \texttt{stringr}}

Cierre de la Parte II. La manipulación de texto es omnipresente en la limpieza de datos: catálogos con códigos pegados a etiquetas, espacios extra, capitalización inconsistente, \textit{paths} de archivos que cambian por trimestre o año. La semana cubre \texttt{stringr} como sistema unificado para operaciones sobre \textit{strings}, las expresiones regulares como lenguaje para describir patrones, \texttt{glue} para construcción dinámica y \texttt{logger} para documentar el flujo del pipeline.

\textbf{Subtemas:}

\begin{itemize}[leftmargin=*, itemsep=2pt]
    \item \textbf{Strings en \textit{base R}.} \texttt{paste()}/\texttt{paste0()}, \texttt{nchar()}, \texttt{tolower()}/\texttt{toupper()}, \texttt{substr()}. Suficientes para tareas triviales; insuficientes para limpieza sistemática de datos.
    \item \textbf{\texttt{stringr}: vocabulario consistente.} El patrón \texttt{str\_*()} ofrece un sistema análogo al de \texttt{dplyr}. Funciones núcleo: detección (\texttt{str\_detect}), extracción (\texttt{str\_extract}, \texttt{str\_match}), reemplazo (\texttt{str\_replace}, \texttt{str\_replace\_all}), división (\texttt{str\_split}), recortado (\texttt{str\_trim}, \texttt{str\_squish}), \textit{padding} (\texttt{str\_pad}), capitalización (\texttt{str\_to\_lower}, \texttt{str\_to\_upper}, \texttt{str\_to\_title}), longitud (\texttt{str\_length}).
    \item \textbf{Expresiones regulares.} Caracteres literales y metacaracteres. Clases (\texttt{[abc]}, \texttt{[a-z]}, \texttt{\textbackslash d}, \texttt{\textbackslash s}, \texttt{\textbackslash w}). Cuantificadores (\texttt{*}, \texttt{+}, \texttt{?}, \texttt{\{n,m\}}). \textit{Anchors} (\texttt{\textasciicircum}, \texttt{\$}). Grupos de captura (\texttt{()}) y referencias hacia atrás (\texttt{\textbackslash 1}). Aplicaciones: validación de formato (CURP, RFC) y extracción de \textit{tokens}.
    \item \textbf{Construcción dinámica de \textit{strings}: \texttt{glue}.} Interpolación de variables (\texttt{glue("Hola \{nombre\}")}), iteración natural sobre vectores, formateo \textit{inline}. Comparación con \texttt{paste0()} y \texttt{sprintf()}. Aplicación a generación de \textit{paths}, mensajes de \textit{log} y plantillas de reporte.
    \item \textbf{Documentación del flujo: \texttt{logger}.} Mensajes estructurados como complemento al \texttt{print()} ocasional: número de filas leídas, descartadas y tiempo por paso. Funciones de \textit{base R} para emitir mensajes (\texttt{message()}, \texttt{warning()}, \texttt{stop()}) y el sistema de niveles del paquete \texttt{logger} (\texttt{log\_info}, \texttt{log\_warn}, \texttt{log\_error}, \texttt{log\_debug}), \textit{timestamps} automáticos y \textit{appenders} configurables. Patrón típico del curso: cada función del pipeline emite \texttt{log\_info()} al inicio y al cierre, dejando un rastro auditable en \texttt{output/log.txt}.
    \item \textbf{Texto como dato: distancias y vectorización.} Aproximación a las cadenas como objeto matemático. El paquete \texttt{stringdist} para distancias entre cadenas (Levenshtein, Jaro-Winkler, Hamming), aplicable a \textit{fuzzy matching} y deduplicación (e.g.\ cruce de nombres de municipios con variaciones ortográficas). Mención breve a \texttt{tidytext} para tokenización y análisis de frecuencias, y a la idea de \textit{word embeddings}.
\end{itemize}

\medskip
\noindent\textbf{Objetivos de aprendizaje}
\begin{enumerate}[leftmargin=*, itemsep=1pt]
    \item Manipular vectores de \textit{strings} con \texttt{stringr}, distinguiendo cuándo basta una operación de \textit{base R} y cuándo se requiere el sistema.
    \item Leer y escribir expresiones regulares para detectar, extraer y reemplazar patrones en datos textuales.
    \item Construir \textit{strings} dinámicos con \texttt{glue}, distinguiéndolos de las alternativas de \textit{base R}.
    \item Limpiar inconsistencias textuales típicas en datos crudos combinando \texttt{stringr} con \textit{regex}.
    \item Documentar el flujo del pipeline con \texttt{logger}, emitiendo mensajes auditables a archivo y consola.
    \item Reconocer cuándo el texto debe tratarse como secuencia de caracteres frente a cuándo conviene tratarlo como objeto matemático.
\end{enumerate}

\medskip
\noindent\textbf{Lecturas}
\begin{itemize}[leftmargin=*, itemsep=1pt]
    \item \textit{R for Data Science} (2.\textsuperscript{a} ed.), Cap.\ 14 \textit{Strings} completo \cite{wickham2023r4ds}.
    \item \textit{R for Data Science} (2.\textsuperscript{a} ed.), Cap.\ 15 \textit{Regular expressions} completo \cite{wickham2023r4ds}.
    \item \textit{Text Mining with R: A Tidy Approach} (Silge \& Robinson), lectura opcional \cite{silge-tidytext}.
\end{itemize}

\medskip
\noindent\textbf{Referencia rápida:} \cheatsheet{strings}{String manipulation with \texttt{stringr}} (Posit); \cheatsheet{regex}{Regular Expressions} (Posit).

% ------------------------------------------------------------------------------
\parteCurso{Parte III --- Programación}
% ------------------------------------------------------------------------------

\semana{8}{Control de flujo y vectorización}

Apertura de la Parte III. La semana cubre la sintaxis de control de flujo en R ---condicionales y \textit{loops}--- y la vectorización como alternativa por defecto del lenguaje. La meta es que el estudiante reconozca cuándo un \textit{loop} es genuinamente necesario y cuándo es un sustituto inadecuado de la operación vectorizada. La semana prepara para los temas siguientes: funciones (S9) y programación funcional con \texttt{purrr} (S10).

\textbf{Subtemas:}

\begin{itemize}[leftmargin=*, itemsep=2pt]
    \item \textbf{Condicionales: \texttt{if}, \texttt{else}, \texttt{switch}.} Sintaxis básica. \texttt{if} opera sobre escalares; aplicado a un vector emite \textit{warning} o evalúa solo el primer elemento. \texttt{switch()} para ramas múltiples mutuamente excluyentes. Distinción: \texttt{if} controla flujo de programa; \texttt{if\_else()} y \texttt{case\_when()} (Semana 6) construyen variables vectorizadas dentro de \texttt{mutate()}.
    \item \textbf{\textit{Loops}: \texttt{for} y \texttt{while}.} Sintaxis e iteración. Iteración segura sobre secuencias con \texttt{seq\_along()} y \texttt{seq\_len()}. Iteración sobre listas. Pre-asignación de la salida como buena práctica; consecuencias del crecimiento del vector dentro del \textit{loop}. Control fino con \texttt{break} y \texttt{next}. Escenarios en que el \textit{loop} se justifica: operaciones dependientes del estado previo, \textit{side effects} (lectura y escritura de archivos), ETL secuencial sobre años o trimestres.
    \item \textbf{Vectorización como modo por defecto.} R aplica la mayoría de operaciones aritméticas, lógicas y funciones elemento a elemento sobre vectores. Comparación entre código con \textit{loop} y código vectorizado, en claridad y rendimiento. Reglas de \textit{recycling}. Funciones acumuladoras (\texttt{cumsum}, \texttt{cummax}, \texttt{cumprod}, \texttt{cumany}, \texttt{cumall}).
    \item \textbf{Patrones y contrapatrones.} Crecimiento de vector (\texttt{x <- c(x, ...)} dentro de \textit{loop}) frente a pre-asignación. Iteración por valores frente a iteración por índices, y los errores \textit{off-by-one}. Mutación involuntaria de objetos dentro de funciones (preludio a la Semana 9, donde se formaliza el \textit{scoping}). Criterio operativo del curso: vectorizar cuando sea posible; cuando no, escribir un \textit{loop} legible y con la salida pre-asignada.
\end{itemize}

\medskip
\noindent\textbf{Objetivos de aprendizaje}
\begin{enumerate}[leftmargin=*, itemsep=1pt]
    \item Escribir condicionales (\texttt{if}/\texttt{else}/\texttt{switch}) y \textit{loops} (\texttt{for}/\texttt{while}) con sintaxis correcta.
    \item Identificar cuándo una operación se puede vectorizar y reescribirla sin \textit{loop} cuando proceda.
    \item Pre-asignar la salida de un \textit{loop} y elegir entre iteración por valores e iteración por índices.
    \item Reconocer contrapatrones frecuentes de iteración y refactorizarlos hacia código vectorizado o \textit{loops} bien construidos.
\end{enumerate}

\medskip
\noindent\textbf{Lecturas}
\begin{itemize}[leftmargin=*, itemsep=1pt]
    \item \textit{Advanced R} (2.\textsuperscript{a} ed.), Cap.\ 5 \textit{Control flow} completo \cite{wickham2019advr}.
\end{itemize}

\semana{9}{Funciones: diseño, scoping y debugging}

La semana marca la transición del estudiante de consumidor de código a autor de código. Las funciones son la unidad fundamental de reuso. La semana cubre la anatomía de una función, el modelo de \textit{scoping} con el que R resuelve nombres, las herramientas para producir y capturar errores, las técnicas para diagnosticarlos cuando ocurren, y el operador \textit{embrace} \texttt{\{\{~\}\}}, que cierra la cuestión planteada en la Semana 5 sobre la interacción entre funciones de usuario y \texttt{dplyr}.

\textbf{Subtemas:}

\begin{itemize}[leftmargin=*, itemsep=2pt]
    \item \textbf{Anatomía y diseño de una función.} Sintaxis (\texttt{function(args) \{ body \}}), argumentos con valor por defecto, argumentos por nombre frente a posicionales, valor de retorno (último valor evaluado o \texttt{return()} explícito). Convención: una función tiene un solo propósito y un nombre que lo expresa. Regla operativa: la repetición de un mismo patrón en tres ocasiones justifica su encapsulamiento en función.
    \item \textbf{\textit{Scoping} léxico.} Resolución de nombres en R: ambiente local, padre y global. Razón por la cual una función no contamina el ambiente global por defecto. Variables libres y \textit{closures}. Distinción entre funciones puras (entrada $\to$ salida, sin efectos colaterales) y funciones con \textit{side effects}. Convención del curso: preferir funciones puras siempre que sea posible.
    \item \textbf{Manejo de errores.} \texttt{stop()} para errores fatales, \texttt{warning()} para advertencias, \texttt{message()} para información. Validación temprana de argumentos con \texttt{stopifnot()} o validaciones manuales con mensajes claros. Captura de errores con \texttt{tryCatch()} para flujos que deben continuar pese a fallas (caso ETL: ante el fallo de un trimestre, proseguir con los siguientes y reportar al cierre).
    \item \textbf{\textit{Debugging}.} \texttt{traceback()} para inspeccionar el \textit{call stack} tras un error. \texttt{browser()} como pausa interactiva en cualquier punto del código. \texttt{debug()} y \texttt{undebug()} para activar el modo paso a paso sobre una función. Estrategia general: lectura del error y del \texttt{traceback}; \textit{print debugging} con \texttt{print()} o \texttt{cat()}; \textit{debugger} interactivo cuando los pasos anteriores resultan insuficientes.
    \item \textbf{Funciones que reciben nombres de columnas: NSE en \texttt{dplyr}.} Razón por la que \texttt{mi\_funcion(df, edad)} falla cuando internamente llama \texttt{filter(df, edad > 18)}. El operador \textit{embrace} \texttt{\{\{~var~\}\}} reenvía un argumento al ámbito de evaluación de \texttt{dplyr}. Operador \texttt{:=} para asignar nombres dinámicos en \texttt{mutate()}. Aplicación: funciones reutilizables sobre cualquier columna del dataset.
\end{itemize}

\medskip
\noindent\textbf{Objetivos de aprendizaje}
\begin{enumerate}[leftmargin=*, itemsep=1pt]
    \item Diseñar funciones de un solo propósito con argumentos definidos y valor de retorno explícito.
    \item Explicar la resolución de nombres dentro de una función y anticipar problemas de \textit{scoping}.
    \item Validar argumentos y producir errores informativos; capturar fallas en flujos secuenciales con \texttt{tryCatch()}.
    \item Diagnosticar errores con \texttt{traceback()} y \texttt{browser()}, distinguiendo cuándo es apropiada cada técnica.
    \item Escribir funciones que reciban nombres de columnas como argumentos mediante el operador \texttt{\{\{~\}\}}.
\end{enumerate}

\medskip
\noindent\textbf{Lecturas}
\begin{itemize}[leftmargin=*, itemsep=1pt]
    \item \textit{Advanced R} (2.\textsuperscript{a} ed.), Cap.\ 6 \textit{Functions} completo \cite{wickham2019advr}.
    \item \textit{Advanced R} (2.\textsuperscript{a} ed.), Cap.\ 8 \textit{Conditions}, §8.2 \textit{Signalling conditions} y §8.4 \textit{Handling conditions} \cite{wickham2019advr}.
    \item \textit{Advanced R} (2.\textsuperscript{a} ed.), Cap.\ 22 \textit{Debugging} completo \cite{wickham2019advr}.
    \item \textit{R for Data Science} (2.\textsuperscript{a} ed.), Cap.\ 25 \textit{Functions}, §25.3 \textit{Data frame functions} \cite{wickham2023r4ds}.
\end{itemize}

\semana{10}{Programación funcional: \texttt{purrr} y \texttt{apply}}

Cierre de la Parte III. La semana aborda la programación funcional como paradigma: tratar funciones como objetos de datos para iterar sobre colecciones sin escribir \textit{loops}. La semana continúa la lógica de las dos anteriores: donde la Semana 8 enseña a evitar \textit{loops} sobre operaciones \textit{built-in}, la Semana 10 enseña a evitar \textit{loops} sobre funciones propias.

\textbf{Subtemas:}

\begin{itemize}[leftmargin=*, itemsep=2pt]
    \item \textbf{El paradigma funcional.} Funciones como objetos de primera clase: asignables, pasables como argumentos, retornables. Razones para preferir \textit{functionals} a un \textit{for}: predictibilidad de tipo, composición y código declarativo. Conexión con la Semana 8: si la operación es \textit{built-in}, vectorización; si es una función propia, \textit{functional}.
    \item \textbf{La familia \texttt{apply} de \textit{base R}.} \texttt{apply()} (filas o columnas de una matriz), \texttt{lapply()} (lista $\to$ lista), \texttt{sapply()} (simplifica cuando puede; tipo de retorno impredecible), \texttt{vapply()} (con tipo declarado), \texttt{mapply()} (varios argumentos), \texttt{tapply()} (\textit{split-apply}). Limitación principal: API inconsistente y tipo de retorno no garantizado.
    \item \textbf{\texttt{purrr}: sistema unificado.} El patrón \texttt{map\_*()} con sufijo de tipo: \texttt{map()} (lista), \texttt{map\_dbl()}, \texttt{map\_chr()}, \texttt{map\_lgl()}, \texttt{map\_dfr()}, \texttt{map\_dfc()} (\textit{data frame}). Tipo de retorno predecible por construcción. Variantes para varias entradas: \texttt{map2()}, \texttt{pmap()}. Variantes para \textit{side effects}: \texttt{walk()}, \texttt{walk2()}. Funciones anónimas con \texttt{\textbackslash(x)} (R 4.1+) o \texttt{\textasciitilde{} .x} (estilo \textit{purrr}).
    \item \textbf{Reducción y predicados.} \texttt{reduce()} para colapsar una lista a un valor; \texttt{accumulate()} para conservar pasos intermedios. Predicados sobre colecciones: \texttt{keep()}, \texttt{discard()}, \texttt{every()}, \texttt{some()}, \texttt{none()}, \texttt{detect()}.
    \item \textbf{\texttt{across()}: iteración dentro de \texttt{dplyr}.} Aplicación de una misma función a varias columnas dentro de \texttt{mutate()} o \texttt{summarize()}. \texttt{if\_any()} y \texttt{if\_all()} para condiciones agregadas dentro de \texttt{filter()}. Aplicaciones al ETL: lectura de todos los CSV de una carpeta, ajuste del mismo modelo a múltiples grupos, escritura de N gráficas.
    \item \textbf{Filtros dinámicos en \texttt{dplyr}: \texttt{expr()} y \texttt{!!}.} Continuación del operador \textit{embrace} introducido en la Semana 9. Cuando lo que varía no es un nombre de columna sino la expresión completa de filtrado (entrada de UI Shiny, lista de condiciones generadas en \texttt{map}), se construye la expresión con \texttt{rlang::expr()} o \texttt{rlang::parse\_expr()} y se inyecta con \texttt{!!}. Sintaxis: \texttt{e <- expr(edad > 30 \& sexo == "M")}; \texttt{filter(datos, !!e)}. Para listas de condiciones: \texttt{filter(datos, !!!lista\_de\_exprs)} con el operador \textit{splice}.
    \item \textbf{Paralelización con \texttt{furrr} y \texttt{future}.} \texttt{furrr::future\_map\_*()} reemplaza a \texttt{purrr::map\_*()} con firmas idénticas; la declaración \texttt{future::plan(multisession, workers = N)} distribuye las iteraciones entre N procesos. Limitación del modelo: R es \textit{single-threaded}; la paralelización ocurre vía procesos separados, no hilos, sin memoria compartida. Los objetos globales se serializan y copian a cada \textit{worker}: N \textit{workers} sobre un dataset de M GB requieren aproximadamente N $\times$ M GB de RAM. La paralelización se justifica cuando la tarea es \textit{CPU-bound} y suficientemente grande para amortizar el costo de \textit{spawn} y serialización.
\end{itemize}

\medskip
\noindent\textbf{Objetivos de aprendizaje}
\begin{enumerate}[leftmargin=*, itemsep=1pt]
    \item Reconocer cuándo una iteración corresponde a un patrón funcional (mapeo, reducción, filtrado por predicado) y elegir el \textit{functional} apropiado.
    \item Usar la familia \texttt{map\_*()} de \texttt{purrr} con tipo de retorno predecible, incluyendo variantes para varias entradas.
    \item Sustituir \textit{loops} explícitos sobre archivos, columnas o grupos por \textit{functionals}, distinguiendo cuándo el \textit{loop} sigue siendo la opción adecuada.
    \item Aplicar \texttt{across()} para vectorizar transformaciones sobre múltiples columnas dentro de \texttt{dplyr}, y construir filtros dinámicos con \texttt{expr()} y \texttt{!!}.
    \item Reconocer las diferencias computacionales entre \textit{loop} \texttt{for} y \textit{functionals} de \texttt{purrr} (estabilidad de tipo, manejo de memoria, ámbito de variables, vía a paralelización).
    \item Paralelizar una iteración con \texttt{furrr} y \texttt{future} cuando se justifique, anticipando las limitaciones de memoria del modelo multi-proceso.
\end{enumerate}

\medskip
\noindent\textbf{Lecturas}
\begin{itemize}[leftmargin=*, itemsep=1pt]
    \item \textit{Advanced R} (2.\textsuperscript{a} ed.), Cap.\ 9 \textit{Functionals} completo \cite{wickham2019advr}.
    \item \textit{R for Data Science} (2.\textsuperscript{a} ed.), Cap.\ 26 \textit{Iteration} completo \cite{wickham2023r4ds}.
\end{itemize}

\medskip
\noindent\textbf{Referencia rápida:} \cheatsheet{purrr}{Apply Functions with \texttt{purrr}} (Posit).

% ------------------------------------------------------------------------------
\parteCurso{Parte IV --- Análisis}
% ------------------------------------------------------------------------------

\semana{11}{Modelado estadístico y series de tiempo}

Apertura de la Parte IV. El enfoque es estrictamente programático: especificación de fórmulas, ajuste de modelos, extracción y reporte de resultados. La teoría estadística subyacente se asume conocida; el curso de Econometría es prerrequisito conceptual de esta semana. La meta operativa es trasladar una pregunta sobre un dataset (e.g.\ ¿qué predice X?, ¿hay diferencia entre grupos?, ¿cómo evoluciona Y en el tiempo?) a una tabla \textit{publication-ready} de resultados.

\textbf{Subtemas:}

\begin{itemize}[leftmargin=*, itemsep=2pt]
    \item \textbf{Fórmulas y \texttt{lm()}.} Sintaxis de fórmulas: operadores (\texttt{+}, \texttt{:}, \texttt{*}, \texttt{\textasciicircum}, \texttt{I()}), control del intercepto (\texttt{-1}/\texttt{+0}). Anatomía del objeto \texttt{lm}: \texttt{summary()}, \texttt{coef()}, \texttt{confint()}, \texttt{predict()}, \texttt{residuals()}, \texttt{fitted()}. Construcción incremental con \texttt{update()}. Diagnósticos: \texttt{plot.lm()} (los cuatro paneles canónicos) y \texttt{vif()} para multicolinealidad.
    \item \textbf{Categóricas en modelos.} Como predictores: \textit{factors} convertidos automáticamente a \textit{dummies}, manipulación del nivel de referencia con \texttt{relevel()} o reordenando el \textit{factor}, interpretación de coeficientes. Como variable dependiente: \texttt{glm(family = binomial)} para modelos logit/probit y \texttt{glm(family = poisson)} para conteos. Interpretación operativa de coeficientes (\textit{odds-ratio}, exponenciación).
    \item \textbf{Datos panel y errores estándar robustos.} \texttt{pdata.frame()} con identificador y dimensión temporal. Estimadores: \texttt{'within'} (efectos fijos), \texttt{'random'} (efectos aleatorios), \texttt{'between'}, \texttt{'pooling'}. Pruebas de Hausman y Breusch-Pagan mediante \texttt{phtest()} y \texttt{plmtest()}. Errores estándar robustos a heteroscedasticidad con \texttt{sandwich::vcovHC()} y reporte vía \texttt{lmtest::coeftest()}, aplicables tanto a \texttt{lm} como a \texttt{plm}.
    \item \textbf{Fechas y series de tiempo (tratamiento mínimo).} Fechas con \texttt{lubridate}: \textit{parsers} (\texttt{ymd()}, \texttt{dmy()}, \texttt{mdy()}), extractores (\texttt{year}, \texttt{month}, \texttt{quarter}, \texttt{wday}), aritmética (\texttt{days()}, \texttt{months()}, \texttt{years()}) y secuencias temporales. Clases para series: \texttt{ts} (regular), \texttt{zoo} (irregular), \texttt{xts} (heredera moderna); conversión entre \texttt{tibble} y series. Operaciones de diagnóstico: \texttt{lag()}, \texttt{diff()}, ACF y PACF visuales.
    \item \textbf{Reportes con \texttt{broom} y \texttt{modelsummary}.} \texttt{broom}: \texttt{tidy()} para coeficientes en \textit{tibble}, \texttt{glance()} para \textit{fit statistics}, \texttt{augment()} para residuales y \textit{fitted} añadidos al \textit{data}. La misma API funciona para \texttt{lm}, \texttt{glm}, \texttt{plm} y otros. \texttt{modelsummary}: comparación de varios modelos lado a lado, personalización (etiquetas, estrellas, decimales), exportación a Word, HTML, Markdown y LaTeX. Comparación con \texttt{stargazer}, herramienta tradicional en economía académica.
\end{itemize}

\medskip
\noindent\textbf{Objetivos de aprendizaje}
\begin{enumerate}[leftmargin=*, itemsep=1pt]
    \item Especificar modelos lineales con la sintaxis de fórmulas de R, manipular el objeto \texttt{lm} y producir diagnósticos básicos.
    \item Incorporar variables categóricas en modelos como predictores y como variable dependiente vía GLM, interpretando los coeficientes en términos operativos.
    \item Estimar modelos de datos panel con \texttt{plm} y reportar errores estándar robustos a heteroscedasticidad con \texttt{sandwich} y \texttt{lmtest}.
    \item Manipular fechas con \texttt{lubridate} y aplicar las clases básicas de series de tiempo (\texttt{ts}, \texttt{zoo}), incluyendo operaciones de diagnóstico (\textit{lag}, diferencias, ACF, PACF).
    \item Producir tablas de modelos \textit{publication-ready} con \texttt{broom} y \texttt{modelsummary}, comparando varios modelos lado a lado.
\end{enumerate}

\medskip
\noindent\textbf{Lecturas}
\begin{itemize}[leftmargin=*, itemsep=1pt]
    \item \textit{An Introduction to R} (manual oficial), Cap.\ 11 \textit{Statistical models in R} \cite{rcoreteam-intro}.
    \item \textit{Panel Data Econometrics in R: The plm Package} (Croissant \& Millo): viñeta principal del paquete \cite{croissant-plm}.
    \item Documentación oficial de \texttt{broom} \cite{robinson-broom} y \texttt{modelsummary} \cite{arelbundock-modelsummary}, secciones \textit{Getting started}.
\end{itemize}

\medskip
\noindent\textbf{Referencia rápida:} \cheatsheet{lubridate}{Dates and times with \texttt{lubridate}} (Posit).

\semana{12}{Datos a escala: SQL, \texttt{DBI} y \texttt{dbplyr}}

Cierre de la Parte IV. La semana cubre las herramientas para conectar R con bases de datos relacionales: SQL como lenguaje, \texttt{DBI} como API de conexión y \texttt{dbplyr} como puente con la sintaxis de \texttt{dplyr}. La meta operativa es mover datos entre archivos, memoria y base de datos, eligiendo con criterio la representación apropiada para cada momento del proyecto.

\textbf{Subtemas:}

\begin{itemize}[leftmargin=*, itemsep=2pt]
    \item \textbf{Motivación y tipología de bases de datos.} Justificación de una base de datos: volumen que excede la memoria, persistencia entre ejecuciones, acceso concurrente, fuentes externas que se exponen vía SQL (CONEVAL, catálogos del INEGI, Banco de México). Comparación: SQLite (archivo, \textit{zero-config}, opción por defecto del curso); PostgreSQL y MySQL (servidores, comunes en producción); duckdb (\textit{zero-config} y columnar, optimizada para \textit{analytics}, particularmente atractiva para encuestas de varios millones de filas). Mención a \texttt{arrow} y \textit{parquet} como alternativa \textit{file-based} para \textit{analytics} a escala.
    \item \textbf{SQL básico.} \texttt{SELECT}, \texttt{FROM}, \texttt{WHERE}, \texttt{ORDER BY}, \texttt{LIMIT}. \texttt{JOIN} (\texttt{INNER}, \texttt{LEFT}, \texttt{RIGHT}, \texttt{FULL}). \texttt{GROUP BY} con funciones agregadoras (\texttt{COUNT}, \texttt{SUM}, \texttt{AVG}, \texttt{MIN}, \texttt{MAX}). Operadores: \texttt{LIKE}, \texttt{IN}, \texttt{BETWEEN}, \texttt{IS NULL}. \textit{Common Table Expressions} con \texttt{WITH}. Mención a \textit{subqueries}. Sentencias de creación: \texttt{CREATE TABLE} e \texttt{INSERT}.
    \item \textbf{\texttt{DBI}: API de conexión.} \texttt{dbConnect()} y \texttt{dbDisconnect()} para abrir y cerrar conexiones; uso de \texttt{tryCatch()} para garantizar el cierre. Operaciones de tabla: \texttt{dbWriteTable()}, \texttt{dbReadTable()}, \texttt{dbListTables()}, \texttt{dbListFields()}. Ejecución de \textit{queries}: \texttt{dbGetQuery()} cuando se espera resultado tabular; \texttt{dbExecute()} cuando no (creación, inserción, actualización). Parametrización segura con \texttt{glue\_sql()} para interpolar valores sin riesgo de inyección SQL.
    \item \textbf{\texttt{dbplyr}: \texttt{dplyr} sobre SQL.} \texttt{tbl(con, 'tabla')} crea un \textit{remote tibble} sobre el cual los verbos de \texttt{dplyr} se traducen automáticamente a SQL y se ejecutan en el servidor. \texttt{show\_query()} expone el SQL generado. \texttt{collect()} materializa el resultado en memoria como \textit{tibble} ordinario. Patrón recomendado: \texttt{filter}, \texttt{group\_by} y \texttt{summarize} en el servidor; \texttt{collect()} al final, cuando el resultado es de tamaño manejable. Limitación: no todas las funciones de R se traducen; las funciones propias no son trasladables.
\end{itemize}

\medskip
\noindent\textbf{Objetivos de aprendizaje}
\begin{enumerate}[leftmargin=*, itemsep=1pt]
    \item Identificar cuándo un proyecto de datos justifica una base de datos relacional y elegir entre las opciones disponibles (SQLite, duckdb, PostgreSQL, MySQL).
    \item Escribir \textit{queries} SQL que cubran las necesidades de análisis comunes: \texttt{SELECT}, \texttt{JOIN}, \texttt{GROUP BY}, agregadoras, \texttt{WITH}.
    \item Conectar R a una base de datos mediante \texttt{DBI}, escribir y leer tablas, y ejecutar \textit{queries} parametrizadas de forma segura.
    \item Utilizar \texttt{dbplyr} para combinar la sintaxis de \texttt{dplyr} con la ejecución en el servidor, e inspeccionar el SQL generado.
\end{enumerate}

\medskip
\noindent\textbf{Lecturas}
\begin{itemize}[leftmargin=*, itemsep=1pt]
    \item \textit{R for Data Science} (2.\textsuperscript{a} ed.), Cap.\ 21 \textit{Databases} completo \cite{wickham2023r4ds}.
    \item Documentación oficial de \texttt{DBI} (sección \textit{Get started}) y \texttt{dbplyr} (\textit{Introduction to dbplyr}, \textit{Verb translation}).
\end{itemize}

% ------------------------------------------------------------------------------
\parteCurso{Parte V --- Productos Finales}
% ------------------------------------------------------------------------------

\semana{13}{Visualización con \texttt{ggplot2}}

Apertura de la Parte V. La semana presenta \texttt{ggplot2} desde la gramática de gráficas: una visualización entendida como composición de capas independientes, cada una con responsabilidades acotadas. La meta es que el estudiante razone por capa y reconozca qué parámetros controla cada una, condición para producir visualizaciones a medida en lugar de copiar plantillas.

\textbf{Subtemas:}

\begin{itemize}[leftmargin=*, itemsep=2pt]
    \item \textbf{Gramática de gráficas y arquitectura por capas.} Origen del enfoque en \textit{Grammar of Graphics} (Wilkinson). Capas que componen una gráfica:
    \begin{itemize}[leftmargin=*, itemsep=1pt, label=$\circ$]
        \item \textbf{Data:} el \textit{tibble} que se grafica.
        \item \textbf{Aesthetics} (\texttt{aes}): mapeo de variable a propiedad visual (\texttt{x}, \texttt{y}, \texttt{color}, \texttt{fill}, \texttt{size}, \texttt{shape}, \texttt{alpha}, \texttt{linetype}).
        \item \textbf{Geoms} (\texttt{geom\_*}): la marca visual (puntos, líneas, barras, áreas).
        \item \textbf{Stats} (\texttt{stat\_*}): transformaciones estadísticas implícitas (\textit{count}, \textit{smooth}, \textit{bin}).
        \item \textbf{Position adjustments:} \texttt{stack}, \texttt{dodge}, \texttt{jitter}, \texttt{fill}.
        \item \textbf{Scales} (\texttt{scale\_*}): traducción de los \textit{aesthetics} a valores visuales.
        \item \textbf{Coordinate systems} (\texttt{coord\_*}): cartesiano, polar, transformaciones de coordenadas.
        \item \textbf{Facets} (\texttt{facet\_*}): partición de la gráfica en sub-paneles.
        \item \textbf{Theme:} tipografía, fondos, ejes, leyendas y demás elementos no asociados a datos.
    \end{itemize}
    \item \textbf{Geometrías y la distinción \textit{setting}/\textit{mapping}.} \textit{Geoms} frecuentes: \texttt{geom\_point}, \texttt{geom\_line}, \texttt{geom\_bar}/\texttt{geom\_col} (\texttt{bar} cuenta, \texttt{col} grafica el valor), \texttt{geom\_histogram}, \texttt{geom\_boxplot}, \texttt{geom\_density}, \texttt{geom\_smooth}, \texttt{geom\_text}/\texttt{geom\_label}, \texttt{geom\_ribbon}, \texttt{geom\_area}, \texttt{geom\_jitter}. Distinción operativa: \texttt{color = grupo} dentro de \texttt{aes()} es \textit{mapping} (la variable controla el color); \texttt{color = 'red'} fuera de \texttt{aes()} es \textit{setting} (color fijo).
    \item \textbf{Scales.} \texttt{scale\_x\_continuous}, \texttt{scale\_x\_discrete}, \texttt{scale\_x\_log10}, \texttt{scale\_x\_date} y sus análogos para \texttt{y}, \texttt{color}, \texttt{fill}, \texttt{size}. Control de \texttt{limits}, \texttt{breaks}, \texttt{labels}, transformaciones. \texttt{scale\_*\_manual()} con vectores nombrados. El paquete \texttt{scales} aporta \textit{formatters} (\texttt{label\_dollar()}, \texttt{label\_percent()}, \texttt{label\_comma()}, \texttt{label\_number\_si()}). Distinción entre \texttt{coord\_cartesian(ylim = ...)} (zoom) y \texttt{scale\_y\_continuous(limits = ...)} (descarta datos).
    \item \textbf{Facets, coordenadas y themes.} \texttt{facet\_wrap} frente a \texttt{facet\_grid}; uso de \texttt{scales = 'free'}. \texttt{coord\_flip()}, \texttt{coord\_polar()}, \texttt{coord\_cartesian()}. Themes predefinidos (\texttt{theme\_minimal}, \texttt{theme\_bw}, \texttt{theme\_classic}, \texttt{theme\_void}) y personalización con \texttt{theme()}. \texttt{theme\_set()} para fijar el tema global del proyecto.
    \item \textbf{Ecosistema complementario y exportación.} Composición y anotación: \texttt{patchwork} (combinación de gráficas con \texttt{+}, \texttt{/}, \texttt{|}), \texttt{ggrepel} (\texttt{geom\_text\_repel} para etiquetas sin solapamiento), \texttt{ggtext} (\textit{markdown}/HTML en títulos y etiquetas), \texttt{gghighlight} (resaltado de grupos). Color: \texttt{ggthemes} (estilos \textit{Economist}, \textit{FiveThirtyEight}, \textit{WSJ}), \texttt{paletteer} (acceso a múltiples paletas). Mención a \texttt{plotly} y \texttt{gganimate}. Exportación con \texttt{ggsave()}: formatos PNG (\textit{raster}, web y Word), JPG (\textit{raster}, calidad inferior) y SVG (vector, escalable sin pérdida). Argumentos: \texttt{width}, \texttt{height}, \texttt{units}, \texttt{dpi}, \texttt{bg = 'white'} para forzar fondo blanco.
\end{itemize}

\medskip
\noindent\textbf{Objetivos de aprendizaje}
\begin{enumerate}[leftmargin=*, itemsep=1pt]
    \item Explicar la arquitectura por capas de \texttt{ggplot2} y el rol de cada una.
    \item Construir una gráfica eligiendo la geometría apropiada y distinguiendo entre \textit{mapping} y \textit{setting}.
    \item Personalizar \textit{scales} para producir ejes profesionales mediante \textit{formatters} del paquete \texttt{scales}.
    \item Aplicar \textit{themes}, particionar con \textit{facets} y elegir la librería complementaria adecuada para cada necesidad.
    \item Exportar gráficas con dimensiones y formato apropiados al medio de destino.
\end{enumerate}

\medskip
\noindent\textbf{Lecturas}
\begin{itemize}[leftmargin=*, itemsep=1pt]
    \item \textit{R for Data Science} (2.\textsuperscript{a} ed.), Cap.\ 9 \textit{Layers} completo \cite{wickham2023r4ds}.
    \item \textit{R for Data Science} (2.\textsuperscript{a} ed.), Cap.\ 11 \textit{Communication} completo \cite{wickham2023r4ds}.
    \item \textit{ggplot2: Elegant Graphics for Data Analysis} (Wickham): referencia comprehensiva \cite{wickham2016ggplot2}.
\end{itemize}

\medskip
\noindent\textbf{Referencia rápida:} \cheatsheet{data-visualization}{Data visualization with \texttt{ggplot2}} (Posit).

\semana{14}{Automatización con \texttt{officer}, \texttt{googledrive} y \texttt{gmailr}}

La semana cubre la transformación de los \textit{outputs} del análisis ---gráficas, tablas, valores clave--- en presentaciones automatizadas, así como el establecimiento de un canal de gestión de archivos mediante Google Drive y de correo mediante Gmail. La intención es sustituir el patrón habitual ---reconstrucción manual de la presentación tras cada actualización de datos--- por un flujo en que la presentación se regenera con un \textit{script} partiendo de un \textit{template} estilizado, y se distribuye al destinatario sin intervención manual. El énfasis del curso se concentra en PPT como medio de entrega.

\textbf{Subtemas:}

\begin{itemize}[leftmargin=*, itemsep=2pt]
    \item \textbf{Motivación y arquitectura del flujo automatizado.} Reconstrucción manual de la presentación tras cada actualización de datos: costo en tiempo, fragilidad y dificultad de auditoría. La automatización opera sobre la fórmula \textit{template} + datos + código = presentación. Patrón general: el pipeline ETL produce \textit{outputs} en \texttt{output/}; un \textit{script} de presentación lee un \texttt{.pptx} pre-estilizado e inyecta los valores actualizados; el resultado se publica vía Drive.
    \item \textbf{\texttt{officer} para PowerPoint.} \texttt{read\_pptx()} para cargar un \textit{template} (preferible a la construcción desde cero). Adición de diapositivas con \texttt{add\_slide()} eligiendo \textit{layouts} del \textit{master} (Título, Contenido, Sección). Inserción de contenido en \textit{placeholders} con nombre mediante \texttt{ph\_with()}: texto, tablas, gráficas \texttt{ggplot2}, imágenes. Iteración sobre múltiples diapositivas para presentaciones con N secciones (e.g.\ un cuadro por estado o por trimestre). Exportación con \texttt{print(doc, target = ...)}. \texttt{officer} también genera Word, fuera del énfasis del curso.
    \item \textbf{Estilización de componentes.} Texto enriquecido: \texttt{fp\_text()} (color, tamaño, negrita, cursiva), \texttt{fp\_par()} (alineación, espaciado), \texttt{fpar()} y \texttt{block\_list()} para composiciones que mezclan formatos en un mismo párrafo. Tablas con \texttt{flextable}: \textit{themes} (\texttt{theme\_vanilla}, \texttt{theme\_box}, personalizado), formato numérico por columna (\texttt{colformat\_double}, \texttt{colformat\_pct}), fusión de celdas, encabezados compuestos, colores por celda. Inserción de gráficas con dimensiones controladas (\texttt{ph\_with(value = gg, location = ph\_location(width, height))}). Convención del curso: todo el estilo reside en el \textit{template} y en el código de \texttt{officer}, no se ajusta en PowerPoint a mano, pues compromete la reproducibilidad.
    \item \textbf{\texttt{googledrive}: autenticación y gestión de archivos.} Autenticación: OAuth interactivo en desarrollo (\texttt{drive\_auth()}); \textit{Service Account} con JSON para automatización (\texttt{drive\_auth(path = "credenciales.json")}). Gestión de archivos: \texttt{drive\_upload()} para publicar reportes; \texttt{drive\_download()} para descargar insumos; \texttt{drive\_ls()} para navegar carpetas; \texttt{drive\_get()} para localizar archivos por nombre o ID; \texttt{drive\_share()} para administrar permisos; \texttt{drive\_mv()} y \texttt{drive\_trash()} para mover y archivar. Patrón habitual: una carpeta de Drive funciona como buzón ---el destinatario deposita insumos ahí, el \textit{script} los descarga, el reporte generado se sube a otra carpeta donde el destinatario lo consume sin contacto con el código.
    \item \textbf{\texttt{gmailr}: correo como canal de automatización.} Autenticación análoga a \texttt{googledrive} sobre la API de Gmail. Envío programático: \texttt{gm\_mime()} para componer un correo (asunto, cuerpo HTML, destinatarios, adjuntos) y \texttt{gm\_send\_message()} para enviarlo. Lectura como buzón de entrada: \texttt{gm\_messages()} (filtros tipo Gmail: \texttt{from:}, \texttt{subject:}, \texttt{has:attachment}), \texttt{gm\_message()}, \texttt{gm\_attachments()}, \texttt{gm\_save\_attachments()}. Patrón: el \textit{script} lee la bandeja con un filtro específico (e.g.\ \texttt{from:proveedor@empresa.com subject:'datos mensuales'}), descarga los adjuntos, los procesa y envía el reporte al destinatario por correo.
\end{itemize}

\medskip
\noindent\textbf{Objetivos de aprendizaje}
\begin{enumerate}[leftmargin=*, itemsep=1pt]
    \item Justificar el patrón \textit{template} + datos + código = presentación frente al flujo manual.
    \item Generar diapositivas de PowerPoint programáticamente a partir de un \textit{template}, insertando texto, tablas y gráficas en \textit{placeholders} con nombre.
    \item Estilizar componentes de una presentación (texto enriquecido, tablas con \texttt{flextable}, imágenes dimensionadas) sin recurrir a edición manual en PowerPoint.
    \item Operar Google Drive desde R para descarga, carga y gestión de archivos, distinguiendo OAuth interactivo de \textit{Service Account}.
    \item Usar \texttt{gmailr} para enviar reportes generados y consumir una bandeja de Gmail como fuente de adjuntos.
    \item Diseñar canales de automatización \textit{low-code} en que carpetas de Drive y bandejas de Gmail operan como buzones de insumos y \textit{outputs}.
\end{enumerate}

\medskip
\noindent\textbf{Lecturas}
\begin{itemize}[leftmargin=*, itemsep=1pt]
    \item Documentación oficial de \texttt{officer}, viñetas \textit{Get started with PowerPoint} y \textit{Layouts and slide manipulation} \cite{gohel-officer}.
    \item Documentación oficial de \texttt{flextable}, secciones \textit{Quick start} y \textit{Theme functions} \cite{gohel-flextable}.
    \item Documentación oficial de \texttt{googledrive}, viñeta \textit{Authentication} y referencia de \texttt{drive\_*} \cite{bryan-googledrive}.
    \item Documentación oficial de \texttt{gmailr}, viñetas \textit{Sending email} y \textit{Working with messages} \cite{hester-gmailr}.
\end{itemize}

\semana{15}{Aplicaciones interactivas con Shiny}

La semana introduce Shiny sin pretensión de exhaustividad. Shiny es una tecnología extensa con uso profesional consolidado; el objetivo del curso no es formar desarrolladores web sino dotar al estudiante de los fundamentos suficientes para construir un tablero funcional (filtros, gráficas interactivas, indicadores numéricos) como producto del proyecto integrador. La semana cubre además los conceptos web subyacentes (HTML, CSS, la capa de conversión a vista, la reactividad) y las limitaciones de la herramienta.

\textbf{Subtemas:}

\begin{itemize}[leftmargin=*, itemsep=2pt]
    \item \textbf{Anatomía de una página web: HTML y CSS.} Modelo cliente-servidor: el navegador solicita, el servidor responde con HTML, CSS y JavaScript. HTML como árbol jerárquico de \textit{tags} (\texttt{div}, \texttt{span}, \texttt{h1}--\texttt{h6}, \texttt{p}, \texttt{a}, \texttt{img}, \texttt{ul}/\texttt{li}); atributos \texttt{id} y \texttt{class}. CSS como sistema de reglas para estilo: selectores por \textit{tag}, \texttt{class} o \texttt{id}; propiedades comunes (\texttt{color}, \texttt{font}, \texttt{background}, \texttt{margin}, \texttt{padding}); estilo \textit{inline} frente a hojas externas. La intención es proveer intuición sobre la pila web, no expertise.
    \item \textbf{La capa R $\to$ HTML/CSS/JS.} Funciones de R que producen HTML: \texttt{tags\$div()}, \texttt{tags\$span()}, \texttt{HTML()}, \texttt{h1()}, \texttt{p()}, \texttt{a()}. Los \textit{inputs} y \textit{outputs} de Shiny son elementos HTML con atributos que el JavaScript del cliente escucha. Bootstrap como librería CSS cargada por defecto. \texttt{htmlwidgets} como puente para integrar librerías JavaScript (highcharts, plotly, leaflet) sin escribir JavaScript.
    \item \textbf{Anatomía mínima de una aplicación Shiny.} Separación UI (vista) y \textit{server} (lógica). \texttt{shinyApp(ui, server)} para definir y \texttt{runApp()} para ejecutar. \textit{Inputs} declarados en UI (\texttt{selectInput}, \texttt{numericInput}, \texttt{sliderInput}, \texttt{dateInput}, \texttt{checkboxInput}, \texttt{textInput}) y accesibles en \textit{server} vía \texttt{input\$id}. \textit{Outputs} declarados en UI (\texttt{plotOutput}, \texttt{tableOutput}, \texttt{textOutput}) y poblados en \textit{server} con \texttt{output\$id <- render*(\{...\})}. Reactividad como modelo de programación: \texttt{reactive()} para valores recalculados al cambiar sus dependencias; \texttt{observe()}/\texttt{observeEvent()} para \textit{side effects}; \texttt{reactiveVal()} para estado mutable explícito; \texttt{isolate()} para suspender la reactividad puntualmente.
    \item \textbf{Estructura de tablero con \texttt{bslib}.} Paquete moderno para \textit{layouts} de tablero. \texttt{page\_sidebar()} para la estructura clásica (filtros a la izquierda, contenido principal a la derecha); \texttt{page\_navbar()} para tableros multi-sección. \texttt{value\_box()} para indicadores numéricos. \texttt{card()} y \texttt{card\_header()} para encapsular secciones. \texttt{layout\_column\_wrap()} y \texttt{layout\_columns()} para \textit{grid responsive}. Gráficas interactivas con \texttt{highcharter} (estética profesional, interactividad nativa) o \texttt{plotly} (integrado al tidyverse vía \texttt{ggplotly()}). Los filtros se conectan a las gráficas mediante reactividad.
    \item \textbf{Limitaciones y criterios de uso.} Cada sesión de Shiny corresponde a un proceso R completo: N usuarios concurrentes implican N veces la memoria. La escalabilidad no es nativa; servir grandes audiencias requiere infraestructura adicional (Posit Connect, Docker). El modelo reactivo se complejiza con interactividad fina. Shiny es apropiado cuando el \textit{output} es R-céntrico (datos, gráficas, modelos), el público es interno o limitado, el tiempo de desarrollo es la restricción dominante y la aplicación tiene propósitos de exploración o reporte interactivo. No es apropiado para escalar a miles de usuarios concurrentes, para interactividad muy granular (React, Vue) o para públicos sin relación con R. Despliegue propuesto: \texttt{shinyapps.io}.
\end{itemize}

\medskip
\noindent\textbf{Objetivos de aprendizaje}
\begin{enumerate}[leftmargin=*, itemsep=1pt]
    \item Explicar la arquitectura cliente-servidor de una página web y los roles de HTML, CSS y JavaScript.
    \item Identificar cómo las funciones de Shiny generan HTML, CSS y JavaScript bajo el capó (\texttt{tags\$*}, Bootstrap, \texttt{htmlwidgets}).
    \item Construir una aplicación Shiny mínima con UI, \textit{server} e interacción reactiva.
    \item Implementar un tablero con filtros, \texttt{value\_box} y gráficas interactivas con \texttt{bslib} junto con \texttt{highcharter} o \texttt{plotly}.
    \item Identificar cuándo Shiny es la herramienta apropiada y cuándo conviene considerar alternativas.
\end{enumerate}

\medskip
\noindent\textbf{Lecturas}
\begin{itemize}[leftmargin=*, itemsep=1pt]
    \item \textit{Mastering Shiny} (Wickham): capítulos introductorios sobre UI básica, \textit{server} y reactividad \cite{wickham2021mastering-shiny}.
    \item Documentación oficial de \texttt{bslib}, viñeta \textit{Dashboards} \cite{bslib-docs}.
    \item Documentación oficial de \texttt{highcharter}, \textit{Get started} \cite{kunst-highcharter}, o \textit{Interactive Web-Based Data Visualization with R, plotly, and shiny} (Sievert), capítulos introductorios \cite{sievert-plotly}.
\end{itemize}

\medskip
\noindent\textbf{Referencia rápida:} \cheatsheet{shiny}{Web Applications with \texttt{shiny}} (Posit).

\semana{16}{Interfaces de modelos de lenguaje}

Última semana del curso. La sesión aborda la incorporación de modelos de lenguaje (LLMs) al pipeline de datos como una herramienta más, no como un servicio externo que interrumpe el flujo. La meta operativa es trasladar al interior del \textit{script} las tareas que habitualmente se delegan a un \textit{chatbot} ---extracción de información estructurada, clasificación de texto, generación de resúmenes, normalización de cadenas---, de modo que sean auditables, reproducibles y ejecutadas sin intervención humana.

\textbf{Subtemas:}

\begin{itemize}[leftmargin=*, itemsep=2pt]
    \item \textbf{Motivación: LLMs integrados al pipeline.} Patrón a sustituir: el analista copia un párrafo no estructurado, lo pega en un \textit{chatbot}, solicita una transformación (extracción, resumen, clasificación), copia la respuesta y la integra al dataset. El procedimiento no es reproducible, no escala, no es auditable y rompe el flujo. La alternativa es invocar al modelo dentro del \textit{script}, como cualquier otra función; cada corrida ejecuta la inferencia y los resultados ingresan al \textit{dataframe} de forma documentada.
    \item \textbf{Dos formas de llamar a una API de LLM.} Forma manual con \texttt{httr2}: construcción del \textit{request} POST al \textit{endpoint} del proveedor (e.g.\ \texttt{/v1/chat/completions}), \textit{headers} con autenticación, \textit{body} JSON, ejecución y parseo. Control total y sin dependencias adicionales, a costa de verbosidad y conocimiento de la API específica. Forma amigable con \texttt{ellmer}: paquete de Posit que abstrae la API en un objeto \texttt{chat} con métodos consistentes y soporte multi-proveedor (Anthropic, OpenAI, Google Gemini, Groq, Ollama local). Mayor portabilidad y soporte para \textit{streaming} y salidas estructuradas, a costa de una capa de abstracción. Convención del curso: \texttt{ellmer} para uso productivo; \texttt{httr2} para control fino o aprendizaje de la mecánica subyacente.
    \item \textbf{Salidas estructuradas.} El uso más valioso de un LLM en un pipeline no es texto libre sino datos estructurados. \texttt{ellmer} soporta esto vía \texttt{chat\$chat\_structured()} con un \textit{schema} declarado con \texttt{type\_object()}, \texttt{type\_string()}, \texttt{type\_number()}, \texttt{type\_array()}. El modelo emite JSON que cumple el \textit{schema}; el parseo produce directamente un \textit{tibble}. Aplicación inmediata: extracción de N campos de un párrafo (acta, noticia, respuesta abierta de encuesta) en una sola pasada lista para \texttt{bind\_rows()}.
    \item \textbf{Casos de uso típicos.} Extracción estructurada (párrafo $\to$ tibble con campos definidos), clasificación (texto $\to$ categoría de una lista cerrada), resumen (texto largo $\to$ resumen corto), normalización (texto no canónico $\to$ canónico), etiquetado (descripción $\to$ lista de \textit{tags}), traducción. Criterio general: tareas en que la transformación de texto a otra forma es rápida para un humano pero tediosa cuando hay N filas, se prestan a automatización con un LLM y un \textit{prompt} adecuado.
    \item \textbf{Extensiones del flujo (panorama).} Modelos locales con \texttt{ollama} para evitar enviar datos a servidores externos en contextos sensibles. OCR local con \texttt{tesseract} para extraer texto de imágenes o PDF escaneados. Detección y supresión de información personal identificable (CURP, RFC, direcciones) mediante regex o LLM. Interfaces conversacionales sobre el pipeline (e.g.\ \texttt{telegram.bot} respondiendo a comandos con análisis). Consideraciones operativas y éticas: costo por \textit{token}, privacidad de los datos enviados a APIs externas, no-determinismo de las respuestas (\texttt{temperature}, \textit{seed}). Regla operativa del curso: si una expresión regular resuelve el problema, una expresión regular es la solución correcta.
\end{itemize}

\medskip
\noindent\textbf{Objetivos de aprendizaje}
\begin{enumerate}[leftmargin=*, itemsep=1pt]
    \item Justificar la integración de LLMs al pipeline de datos frente a su uso como herramienta externa.
    \item Distinguir \texttt{httr2} (manual, transparente) de \texttt{ellmer} (abstracción), eligiendo entre ambos según la situación.
    \item Diseñar y consumir salidas estructuradas con \textit{schemas} para integración directa al \textit{dataframe}.
    \item Aplicar un LLM a tareas características de proyectos de datos: extracción estructurada, clasificación, resumen, normalización, etiquetado.
    \item Reconocer el panorama de extensiones (modelos locales, OCR, PII, interfaces conversacionales) y las consideraciones operativas y éticas asociadas.
\end{enumerate}

\medskip
\noindent\textbf{Lecturas}
\begin{itemize}[leftmargin=*, itemsep=1pt]
    \item Documentación oficial de \texttt{ellmer}, viñetas \textit{Getting started} y \textit{Structured data} \cite{wickham-ellmer}.
    \item Documentación oficial de \texttt{httr2}, viñetas \textit{Getting started} y \textit{Wrapping APIs} \cite{wickham-httr2}.
    \item Documentación del proveedor LLM elegido (Anthropic, OpenAI, Google Gemini, Groq), referencia del \textit{endpoint} de \textit{chat completions} y de \textit{tool use} / \textit{structured output}.
\end{itemize}

% ==============================================================================
\section{Proyecto integrador}
% ==============================================================================

El curso se articula alrededor de un proyecto integrador que cada estudiante construye de forma individual a lo largo del semestre, replicando el flujo completo de un proyecto de datos.

\subsection*{Dataset y restricciones}

Cada estudiante elige una encuesta del INEGI sujeta a tres restricciones:

\begin{enumerate}[leftmargin=*, itemsep=1pt]
    \item Microdatos, no agregados ni tabulados pre-procesados.
    \item Al menos dos módulos relacionables (e.g.\ ENSU con CS/CB/VIV; ENIGH con viviendas/hogares/personas/concentrado), de modo que el ejercicio de \textit{join} entre tablas sea sustantivo.
    \item Al menos dos levantamientos en el tiempo (panel, trimestral o varios años), de modo que el manejo de fechas y la acumulación temporal estén justificados.
\end{enumerate}

Encuestas que satisfacen estas restricciones: ENSU, ENOE, ENIGH, ENVIPE, ENCIG, ENIF, ENADID, MCS-ENIGH, entre otras.

\subsection*{Pipeline}

El proyecto sigue el patrón ETL + análisis + producto:

\begin{enumerate}[leftmargin=*, itemsep=1pt]
    \item \textbf{Extract:} descarga de los archivos de microdatos desde el sitio del INEGI; lectura de los descriptores técnicos.
    \item \textbf{Transform:} corrección de tipos, recodificación de categóricas, manejo de \texttt{NA}, \textit{joins} entre módulos y \textit{reshape} a formato \textit{tidy}.
    \item \textbf{Load:} guardado de los datos transformados en un formato eficiente y reproducible (\texttt{.rds} y, hacia el cierre del semestre, base SQLite local).
    \item \textbf{Análisis:} aplicación de modelos estadísticos apropiados a las preguntas planteadas.
    \item \textbf{Producto:} tablero Shiny con visualizaciones interactivas y reporte periódico automatizado con \texttt{officer}.
\end{enumerate}

\subsection*{Checkpoints}

El proyecto se evalúa por cinco \textit{checkpoints} alineados con las cinco partes del curso:

\begin{center}
\begin{tabular}{cp{2.6cm}p{9cm}}
\toprule
\textbf{\#} & \textbf{Cierre de} & \textbf{Entregable} \\
\midrule
1 & Parte I (Sem.\ 4)   & Repositorio Git con estructura \texttt{files/}, \texttt{output/}, \texttt{docs/}, \texttt{pre/}; justificación de la encuesta elegida (una página); lectura del descriptor técnico; primer \textit{script} de extracción. \\
2 & Parte II (Sem.\ 7)  & ETL completo en \textit{scripts}: \textit{extract} + \textit{transform} + \textit{load} a \texttt{.rds}. Datos \textit{tidy} y validados. \\
3 & Parte III (Sem.\ 10) & Refactor del ETL en funciones reutilizables (\texttt{code/funciones.R}); \textit{master script} con \texttt{source()}. \\
4 & Parte IV (Sem.\ 12) & Análisis estadístico aplicado; \textit{load} conectado a base SQLite local. \\
5 & Parte V (Sem.\ 16) & Tablero Shiny funcional y reporte periódico automatizado con \texttt{officer}. \\
\bottomrule
\end{tabular}
\end{center}

% ==============================================================================
\section{Referencias principales}
% ==============================================================================

Las lecturas específicas por sesión se indican en el temario detallado. Las obras de consulta general son:

\nocite{wickham2023r4ds, wickham2019advr, wickham2016ggplot2, wickham2021mastering-shiny, rcoreteam-intro, bryan-happygit, ushey-renv, silge-tidytext, croissant-plm, robinson-broom, arelbundock-modelsummary, gohel-officer, gohel-flextable, bryan-googledrive, hester-gmailr, bslib-docs, kunst-highcharter, sievert-plotly, wickham-ellmer, wickham-httr2}

\printbibliography[heading=none]

\end{document}
